\section{Contexto del proyecto y tratamientos}

\subsection{Suposiciones}

\begin{itemize}
  \item Se asume que los modelos matemáticos y de control desarrollados representan con suficiente precisión el comportamiento dinámico del convertidor y sus componentes.
  \item Se supone que la plataforma HIL es compatible y permitirá una integración efectiva con el hardware físico del convertidor.
  \item Se considera que los módulos físicos del convertidor están disponibles y cumplen con las especificaciones técnicas necesarias para la implementación y pruebas.
  \item Se asume que las condiciones de operación simuladas y experimentales serán representativas de escenarios reales de funcionamiento en microrredes y sistemas de almacenamiento de energía.
    \item Se considera que la resolución temporal y capacidad de procesamiento de la plataforma HIL son suficientes para replicar con precisión el comportamiento dinámico del sistema real.

\end{itemize}

\subsection{Restricciones}

\begin{itemize}
  \item El desarrollo experimental del proyecto depende de la disponibilidad de los módulos físicos de H-Bridge.
  \item El acceso al laboratorio, equipos y soporte técnico de la universidad Distrital está condicionado a los tiempos, políticas y disponibilidad de los profesores y personal de apoyo.
  \item El cronograma del proyecto puede verse afectado por los trámites administrativos requeridos para el ingreso y utilización de los recursos en la universidad colaboradora.
  \item Las pruebas físicas del convertidor en configuración interleaved estarán limitadas al equipamiento previamente disponible, sin posibilidad de rediseño estructural de hardware.


\end{itemize}

\subsection{Factores de riesgo}

\begin{itemize}
  \item Posibles demoras en la disponibilidad de los módulos o recursos necesarios en la universidad colaboradora, lo cual podría retrasar la etapa de implementación y validación.
  \item Restricciones administrativas o logísticas que dificulten el acceso oportuno a laboratorios o equipamiento.
  \item Incompatibilidad entre la plataforma HIL del proyecto y el diseño físico de los módulos existentes, lo cual podría requerir adaptaciones técnicas imprevistas.
  \item Riesgos asociados al transporte, manipulación o fallas en los módulos físicos que comprometan el desarrollo del prototipo interleaved.
\end{itemize}

